% 第八章 政策与地缘政治分析

功率半导体行业处于政策红利与地缘风险的双向拉扯之中。国内政策端持续加码——大基金三期落地、税收优惠延续、新能源需求政策发力——为行业提供长期成长支撑；与此同时，出口管制升级、国际贸易壁垒、供应链安全等外部风险构成实质性约束。整体判断：政策红利大于地缘风险，国产替代逻辑在外部压力下持续强化，但短期需高度关注出口管制升级和贸易摩擦的阶段性冲击。

本章从国内产业政策、出口管制、国际贸易与供应链安全、新能源政策联动、地缘政治情景五个维度展开系统分析，旨在为投资决策提供政策面的完整评估框架。

\section{国内产业政策环境}

\subsection{大基金三期：产业升级的核心抓手}

国家集成电路产业投资基金三期（以下简称"大基金三期"）于2024年5月注册成立，注册资本3440亿元，存续期15年（2024--2039年），是当前国内半导体产业最重要的政策资金来源。与二期相比，三期资金规模更大、投向更聚焦于上游核心环节，对功率半导体产业的支持力度显著提升。

\begin{exhibitbox}[表 8-1：大基金三期概况与投资结构]
\centering
\small
\begin{tabularx}{\textwidth}{X L{5.5cm} c}
\toprule
\textbf{项目} & \textbf{详情} & \textbf{状态} \\
\midrule
基金规模 & 注册资本3440亿元 & \textcolor{navy}{\textbf{已落地}} \\
资金来源 & 中央财政、地方政府、国有金融机构、产业资本 & \textcolor{navy}{\textbf{已落地}} \\
存续期限 & 15年（2024--2039年） & \textcolor{navy}{\textbf{已落地}} \\
管理人 & 华芯投资管理有限责任公司 & \textcolor{navy}{\textbf{已落地}} \\
\midrule
\rowcolor{lightgrey}
\multicolumn{3}{l}{\textit{投向结构估算}} \\
设备与材料 & 约占40--45\%，重点支持光刻机、刻蚀机、硅片、光刻胶等 & 规划中 \\
芯片设计与制造 & 约占35--40\%，包括功率半导体、模拟芯片、车规芯片 & 规划中 \\
封测与下游应用 & 约占15--20\%，包括先进封装、汽车电子、工业控制 & 规划中 \\
\midrule
投资节奏 & 2025--2027年为投资高峰期，年均约600--800亿元 & 进行中 \\
\bottomrule
\end{tabularx}
\sourcenote{国家集成电路产业投资基金三期设立公告、工信部公开信息、本研究整理}
\end{exhibitbox}

大基金三期对功率半导体行业的影响通过四条路径传导：一是直接股权投资，为行业龙头提供低成本资金，支持扩产和研发；二是加速上游设备材料国产替代，缓解功率半导体企业的扩产瓶颈；三是推动产能扩张，尤其是SiC/GaN等第三代半导体产能；四是政策背书带来估值溢价，提升板块整体估值中枢。

从已落地投资情况看，大基金三期已对多家功率半导体及上游企业进行投资，覆盖设备、材料、设计、制造各环节，其中第三代半导体（SiC/GaN）是重点支持方向，覆盖衬底、外延、器件全链条。2025--2027年是大基金三期的投资高峰期，预计年均投资规模约600--800亿元，将持续为板块带来催化。

\subsection{税收优惠：实际净利润贡献5--15\%}

税收优惠是半导体产业政策体系的重要组成部分，对功率半导体企业的实际盈利具有显著影响。当前行业适用的税收优惠政策体系涵盖企业所得税、增值税、进口关税等多个维度。

\begin{exhibitbox}[表 8-2：主要税收优惠政策及对功率半导体企业影响]
\centering
\small
\begin{tabularx}{\textwidth}{X L{4.5cm} L{2.5cm} c}
\toprule
\textbf{政策类型} & \textbf{优惠内容} & \textbf{适用条件} & \textbf{状态} \\
\midrule
高新技术企业 & 减按15\%税率征收企业所得税（法定25\%） & 核心自主知识产权、研发投入达标 & 已落地 \\
集成电路生产企业 & 28nm及以下10年免税；65nm及以下5免5减半；130nm及以下2免3减半 & 符合生产线宽条件 & 已落地 \\
集成电路设计企业 & 2免3减半（前2年免税，后3年减半征收） & 符合设计企业认定条件 & 已落地 \\
重点软件/集成电路企业 & 减按10\%税率征收 & 国家规划布局内重点企业 & 已落地 \\
研发费用加计扣除 & 科技型中小企业100\%加计扣除；一般企业75\% & 研发费用归集规范 & 已落地 \\
\midrule
增值税留抵退税 & 集成电路重大项目增值税留抵退税 & 列入国家重大项目 & 已落地 \\
软件产品增值税 & 超3\%即征即退 & 软件产品登记 & 已落地 \\
设备进口免税 & 关键设备、零部件、原材料免征进口关税和增值税 & 符合免税目录 & 已落地 \\
\bottomrule
\end{tabularx}
\sourcenote{财政部、税务总局、工信部相关政策文件，本研究整理}
\end{exhibitbox}

根据重点覆盖公司2025年年报数据，税收优惠对净利润的贡献约为5--15\%。其中，华润微、斯达半导等拥有制造产线的企业受益更为显著，实际所得税税率普遍在10--15\%区间，显著低于25\%的法定税率。研发费用加计扣除政策对高研发投入的SiC相关企业（如天岳先进、三安光电）同样具有重要影响，但在亏损阶段实际税盾效应有限。

\subsection{"十五五"规划前瞻：长期成长空间再确认}

当前"十五五"（2026--2030年）规划尚在编制阶段，但已释放的政策信号显示，半导体产业和新能源产业的战略地位将进一步提升。对功率半导体行业影响最为显著的政策方向包括：

\textbf{电网投资加码。}"十五五"期间电网固定资产投资预计达4万亿元，较"十四五"增长约40\%。特高压建设提速和智能电网升级，将对高压IGBT、晶闸管等产品形成长期需求支撑。时代电气等具备高压大功率器件能力的厂商将显著受益。

\textbf{第三代半导体上升为国家战略。}SiC/GaN等宽禁带半导体有望在"十五五"期间上升为国家战略级产业，车规级芯片国产化率目标将进一步提高，设备材料自主可控成为核心任务。这将从研发投入、产能建设、应用推广等多个维度全面支持第三代半导体产业发展。

\textbf{新能源装机目标上调。}光伏装机2030年目标有望上调至1500GW以上，新型储能2030年目标有望超过150GW，新能源车渗透率目标持续提升。下游需求政策的加码将直接拉动功率半导体市场空间扩容。

\begin{keyinsight}[国内政策环境评估]
国内产业政策对功率半导体行业形成全方位、多层次的支持体系。大基金三期提供资金支持，税收优惠直接增厚利润，"十五五"规划打开长期成长空间。三重政策红利叠加，构成行业长期成长的确定性基石。投资上，建议重点关注政策敏感度高、直接受益于国产替代和第三代半导体政策的标的。
\end{keyinsight}

\section{出口管制与技术封锁影响}

\subsection{管制体系演进与当前格局}

美国对华半导体管制经历了从"点"到"面"的逐步升级过程。2018年《出口管制改革法案》奠定法律基础，2022年《芯片与科学法案》和BIS出口管制规则大幅收紧，2023年AI芯片管制升级，2024年以来半导体设备管制持续扩围。当前管制体系主要通过三个维度实施：实体清单、出口管制分类编码（ECCN）、外国直接产品规则（FDPR）。

\begin{exhibitbox}[表 8-3：美国对华半导体管制政策时间线]
\centering
\small
\begin{tabularx}{\textwidth}{X L{3.5cm} L{4.2cm} L{3cm}}
\toprule
\textbf{时间} & \textbf{政策事件} & \textbf{核心内容} & \textbf{对功率半导体影响} \\
\midrule
2018.08 & 《出口管制改革法案》(ECRA) & 扩大出口管制范围，建立新兴技术管制机制 & 间接影响，奠定法律基础 \\
2020.09 & 华为管制升级 & 限制使用美国技术/设备生产的芯片供应华为 & 对华为供应链企业有影响 \\
2022.08 & 《芯片与科学法案》 & 527亿美元补贴美国半导体产业，限制获补贴企业在华扩产 & 限制外资在华功率半导体产能扩张 \\
2022.10 & BIS更新出口管制规则 & 限制先进逻辑芯片、计算芯片及相关设备出口 & 直接影响有限，设备端影响较大 \\
2023.10 & AI芯片出口管制升级 & 限制H100/A100等高端AI芯片出口 & 对电源管理芯片有间接影响 \\
2024年 & 半导体设备管制扩围 & 更多刻蚀、薄膜沉积、离子注入设备纳入管制 & 对扩产有直接影响 \\
2025年 & 管制进一步收紧传闻 & 传闻扩大成熟制程设备管制、增加第三代半导体管制 & 若落地将显著影响行业 \\
\bottomrule
\end{tabularx}
\sourcenote{BIS官网、美国白宫公告、公开信息整理}
\end{exhibitbox}

\subsection{功率半导体在管制体系中的位置}

总体判断：器件端直接管制有限，设备端间接影响显著。功率半导体器件本身（IGBT、MOSFET、SiC二极管等）大部分未被直接纳入出口管制清单，消费级、工业级产品贸易基本正常。车规级产品未被专门管制，但部分高端车规级SiC器件可能因技术参数触及管制阈值。

真正构成实质性约束的是上游设备和材料环节。高端光刻机、离子注入机、外延设备、刻蚀机等均受管制，是功率半导体扩产的主要瓶颈。SiC相关设备与材料（高温离子注入、SiC晶体生长炉、高温退火设备等）的管制呈收紧趋势。

\begin{exhibitbox}[表 8-4：功率半导体各环节管制程度评估]
\centering
\small
\begin{tabularx}{\textwidth}{X c L{7cm}}
\toprule
\textbf{环节} & \textbf{管制程度} & \textbf{具体说明} \\
\midrule
功率器件成品 & \textcolor{riskgreen}{低} & 大部分器件未被直接纳入出口管制清单，消费级、工业级贸易基本正常 \\
车规级功率器件 & \textcolor{riskgreen}{低--中} & 车规级产品未被专门管制，部分高端SiC器件可能触及技术参数阈值 \\
功率半导体设备 & \textcolor{riskred}{高} & 高端光刻机、离子注入机、外延设备、刻蚀机等均受管制，是扩产主要瓶颈 \\
SiC衬底/外延片 & \textcolor{riskamber}{中} & 高端SiC衬底（6英寸以上、低缺陷密度）部分受管制，普通工业级相对宽松 \\
EDA软件 & \textcolor{riskred}{高} & 先进制程EDA受严格管制，功率半导体特色工艺EDA受影响相对较小 \\
关键材料 & \textcolor{riskamber}{中} & 高端光刻胶、电子特气等部分受管制 \\
\bottomrule
\end{tabularx}
\sourcenote{BIS出口管制清单、行业调研、本研究评估}
\end{exhibitbox}

\subsection{设备进口限制的具体影响}

设备是功率半导体扩产的核心瓶颈。受管制影响最显著的设备类型包括：DUV光刻机（193nm沉浸式及以上）、高能高电流离子注入机、高端外延设备、高深宽比刻蚀机、ALD/高端PVD/CVD等。其中，离子注入机和外延设备对功率半导体制造尤为关键——IGBT和SiC器件对离子注入工艺要求高，功率器件对外延层质量要求严格。

\begin{exhibitbox}[表 8-5：主要功率半导体企业受设备管制影响评估]
\centering
\small
\begin{tabularx}{\textwidth}{l X c c L{3.5cm}}
\toprule
\textbf{公司} & \textbf{扩产计划} & \textbf{设备依赖度} & \textbf{受影响程度} & \textbf{应对措施} \\
\midrule
华润微 & 12英寸产线扩产 & 高 & \textcolor{riskred}{高} & 提前锁定设备、国产替代 \\
士兰微 & 12英寸模拟线、8英寸SiC & 高 & \textcolor{riskred}{高} & 多元设备来源、国产验证 \\
时代电气 & SiC三期、IGBT扩产 & 中高 & \textcolor{riskamber}{中高} & 提前备货、国产设备验证 \\
三安光电 & SiC/GaN产能扩张 & 高 & \textcolor{riskred}{高} & 设备提前布局、自研部分设备 \\
扬杰科技 & 8英寸线、SiC扩产 & 中高 & \textcolor{riskamber}{中高} & 国产设备导入、海外产能 \\
天岳先进 & SiC衬底扩产 & 中高 & \textcolor{riskamber}{中高} & 长晶炉部分自研、国产设备 \\
斯达半导 & 模块封装扩产、晶圆产能 & 中 & \textcolor{riskamber}{中} & 代工为主，设备依赖度较低 \\
\bottomrule
\end{tabularx}
\sourcenote{各公司公告、IR交流、本研究评估}
\end{exhibitbox}

从风险暴露程度看，IDM模式且有大规模12英寸产线扩产计划的企业（华润微、士兰微）受设备管制影响最大；SiC产能扩张企业（三安光电、时代电气、天岳先进）受专用设备限制影响显著；Fabless模式企业（新洁能、东微半导）通过代工模式分散设备风险，受影响相对较小。

\subsection{材料进口风险}

关键材料的进口依赖是另一重要风险维度。高端硅片（12英寸约70\%进口依赖）、高端光刻胶（约80\%进口依赖）、SiC衬底（高端约50\%进口依赖）、电子特气（高端约60\%进口依赖）是最为薄弱的环节。

出口管制升级对材料端的传导路径有三：一是关键材料供应受限导致晶圆厂产能利用率下降；二是材料价格上涨推升制造成本、压缩毛利率；三是倒逼国产材料验证加速，长期利好国产材料企业。从已观察到的影响看，高端6英寸导电型SiC衬底进口周期已延长2--3个月，价格波动加大；ArF光刻胶供应紧张，价格上涨约15--20\%。

\begin{keyinsight}[出口管制影响评估]
出口管制对功率半导体行业的影响呈现"短期承压、长期利好"的特征。短期看，设备和材料进口受限拖累扩产节奏、推升生产成本；中长期看，管制倒逼自主可控，加速设备材料国产化进程，强化国产替代的确定性。投资上，需区分短期情绪冲击和长期逻辑变化，在管制升级带来的估值回调中布局具备长期竞争力的龙头企业。
\end{keyinsight}

\section{国际贸易壁垒与供应链安全}

\subsection{欧盟CBAM：当前影响有限，长期需警惕扩围风险}

欧盟碳边境调节机制（CBAM）于2026年正式实施，首批覆盖钢铁、铝、水泥、化肥、电力、氢气六大行业。半导体目前未被纳入CBAM覆盖范围，对功率半导体行业的直接影响有限。

\begin{exhibitbox}[表 8-6：主要功率半导体企业欧洲业务占比与CBAM影响]
\centering
\small
\begin{tabularx}{\textwidth}{X c c c}
\toprule
\textbf{公司} & \textbf{欧洲收入占比（估算）} & \textbf{CBAM当前影响} & \textbf{未来风险} \\
\midrule
扬杰科技 & 约3--4\% & \textcolor{riskgreen}{低} & \textcolor{riskamber}{中} \\
时代电气 & 约2--3\% & \textcolor{riskgreen}{低} & \textcolor{riskgreen}{低--中} \\
斯达半导 & 约2--3\% & \textcolor{riskgreen}{低} & \textcolor{riskgreen}{低--中} \\
华润微 & 约1--2\% & \textcolor{riskgreen}{低} & \textcolor{riskgreen}{低} \\
士兰微 & 约2--3\% & \textcolor{riskgreen}{低} & \textcolor{riskgreen}{低--中} \\
天岳先进 & 海外占比高（含欧洲客户） & \textcolor{riskgreen}{低--中} & \textcolor{riskamber}{中高} \\
\bottomrule
\end{tabularx}
\sourcenote{各公司年报、IR交流、本研究测算}
\end{exhibitbox}

需警惕的是长期扩围风险。欧盟已明确表示将逐步扩大CBAM覆盖范围，2030年前可能纳入更多行业。若半导体被纳入，高能耗的晶圆制造和SiC衬底生产环节将受较大影响。从企业层面看，天岳先进等海外收入占比较高的SiC衬底企业面临的潜在风险相对更大。

\subsection{美国关税政策与供应链重构}

中美贸易摩擦以来，部分半导体产品被加征25\%关税。中国功率半导体对美出口规模相对有限，且以中低端产品为主，直接影响可控。但关税政策的不确定性推动国内厂商加速东南亚产能布局，以规避贸易壁垒、分散供应链风险。

目前国内功率半导体企业的东南亚布局仍处于早期阶段。扬杰科技越南工厂已实现二期小信号产线通线满产，定位封装测试和部分制造环节；天岳先进马来西亚工厂处于建设中，定位SiC衬底制造；时代电气印尼电驱基地已投产，但主要是汽车电驱系统而非半导体制造；多数企业仍处于前期调研或规划阶段。

我们判断，功率半导体制造环节因技术壁垒高、产业链配套要求高，大规模向东南亚转移的可能性较低。封装测试等后端环节可能逐步外迁，但核心制造能力仍将保留在国内。未来更可能形成"国内+东南亚"双基地格局——高端制造留在国内，中低端封装测试和部分产能布局东南亚以服务海外客户。

\subsection{供应链安全评估}

供应链安全是当前功率半导体行业的核心议题之一。从全产业链视角看，上游设备和材料环节最为薄弱，自主可控程度最低；设计环节中低端已基本自主，高端仍有差距；制造环节特色工艺逐步突破，车规级等高端工艺仍在追赶；封测环节自主可控程度相对较高。

\begin{exhibitbox}[表 8-7：功率半导体供应链薄弱环节评估]
\centering
\small
\begin{tabularx}{\textwidth}{X L{3cm} c L{3.5cm} c}
\toprule
\textbf{环节} & \textbf{细分领域} & \textbf{自主可控程度} & \textbf{薄弱点} & \textbf{风险等级} \\
\midrule
材料 & 高端硅片（12英寸） & 30--40\% & 大尺寸、高纯度硅片 & \textcolor{riskred}{高} \\
材料 & SiC衬底（高端） & 20--30\% & 6英寸以上、低缺陷密度 & \textcolor{riskred}{高} \\
材料 & 光刻胶（ArF/EUV） & <20\% & 高端光刻胶几乎完全依赖进口 & \textcolor{riskred}{很高} \\
设备 & 光刻机（DUV/EUV） & <10\% & 高端光刻机几乎完全依赖进口 & \textcolor{riskred}{很高} \\
设备 & 离子注入机 & 20--30\% & 高能、高电流离子注入 & \textcolor{riskred}{高} \\
设备 & 外延设备（高端） & 20--30\% & SiC外延、高端硅外延 & \textcolor{riskred}{高} \\
设备 & 刻蚀/薄膜沉积 & 30--50\% & 中低端已突破，高端仍依赖 & \textcolor{riskamber}{中高} \\
设计 & EDA工具（高端） & <30\% & 数字全流程、先进节点 & \textcolor{riskred}{高} \\
制造 & 高端工艺 & 40--50\% & 车规级、高可靠性工艺 & \textcolor{riskamber}{中高} \\
封测 & 先进封装 & 50--60\% & 高端封装基板依赖进口 & \textcolor{riskamber}{中} \\
\bottomrule
\end{tabularx}
\sourcenote{中国半导体行业协会、券商研报汇总、本研究评估}
\end{exhibitbox}

令人欣慰的是，在政策推动和市场需求拉动下，国内设备材料企业正加速突破。刻蚀机（中微公司、北方华创）、PVD/CVD（北方华创、拓荆科技）、CMP设备（华海清科）、清洗设备（盛美上海）等已实现中高端突破并进入量产；i线/g线光刻机、离子注入机、外延设备等也在加快验证和国产替代进程。大基金三期将设备材料作为第一投资重点（约40--45\%），将进一步加速这一进程。

\subsection{供应链风险矩阵与缓释措施}

为有效管理供应链风险，我们构建了风险评估矩阵，从发生概率和影响程度两个维度对主要风险事件进行评估，并提出相应的缓释措施。

\begin{exhibitbox}[表 8-8：功率半导体供应链风险矩阵]
\centering
\small
\begin{tabularx}{\textwidth}{X c c c L{3.2cm}}
\toprule
\textbf{风险事件} & \textbf{发生概率} & \textbf{影响程度} & \textbf{风险等级} & \textbf{预警信号} \\
\midrule
美国扩大成熟制程设备管制 & 中（30--40\%） & 很高 & \textcolor{riskred}{\textbf{高}} & BIS新规发布、政策传闻 \\
关键设备交付延期 & 高（50--60\%） & 中 & \textcolor{riskamber}{\textbf{中高}} & 设备厂交期延长、产能紧张 \\
SiC衬底进口受限 & 中（25--35\%） & 高 & \textcolor{riskamber}{\textbf{中高}} & Wolfspeed等出口限制 \\
海外厂商技术封锁升级 & 中（30--40\%） & 中高 & \textcolor{riskamber}{\textbf{中高}} & 专利诉讼、技术授权限制 \\
国内设备材料验证不及预期 & 中高（40--50\%） & 中 & \textcolor{riskamber}{\textbf{中高}} & 国产设备良率、稳定性数据 \\
高端光刻胶断供 & 低（15--20\%） & 高 & \textcolor{riskamber}{\textbf{中高}} & 日本出口管制、产能事故 \\
电子特气供应中断 & 低（10--15\%） & 中高 & \textcolor{riskamber}{\textbf{中}} & 地缘冲突、产能事故 \\
封装基板供应紧张 & 中（30--40\%） & 中 & \textcolor{riskamber}{\textbf{中}} & 基板厂产能利用率、价格走势 \\
\bottomrule
\end{tabularx}
\sourcenote{本研究团队基于公开信息和行业调研构建}
\end{exhibitbox}

针对上述风险，行业普遍采用的缓释措施包括：关键设备和材料保持6--12个月安全库存、建立2--3家合格供应商分散供应风险、加速国产设备材料验证导入、推动上游配套企业就近布局、加大研发投入突破关键核心技术等。从A股公司的风险暴露看，士兰微和华润微因12英寸产线扩产对高端设备依赖度高而处于高风险暴露区间，斯达半导和扬杰基因代工+自有产线模式风险相对分散，新洁能和东微半导等Fabless公司风险暴露最低。

\section{新能源汽车/光伏政策联动效应}

\subsection{新能源汽车政策：购置税减免延续至2027年}

新能源汽车是功率半导体第一大下游（约占35--40\%需求），新能源汽车政策对功率半导体需求具有决定性影响。当前核心政策是新能源汽车车辆购置税减免——根据财政部、税务总局、工信部2023年第10号公告，2024--2025年全额免征车辆购置税，2026--2027年减半征收，每辆新能源乘用车免税额分别不超过3万元和1.5万元。

购置税减免政策对新能源车销量的拉动效应约为5--10个百分点，通过"销量增长+单车价值量提升"双轮驱动功率半导体需求。2025年中国新能源车销量预计约1200--1300万辆，单车功率半导体价值量约1500--2500元，对应车规功率半导体市场规模约200--250亿元，到2030年有望增长至450--550亿元。

需关注的风险点是2027年底购置税减免到期后的政策衔接。若政策退坡力度过大，可能对新能源车需求造成阶段性冲击，进而影响车规功率半导体需求节奏。

\subsection{光伏与储能政策：装机目标持续上调}

光伏和储能是功率半导体的重要增长极。在政策持续推动下，光伏和储能装机目标不断上调，实际装机持续超预期。

\begin{exhibitbox}[表 8-9：光伏与储能核心政策及装机目标]
\centering
\small
\begin{tabularx}{\textwidth}{X L{2cm} L{3cm} L{3cm} c}
\toprule
\textbf{政策文件} & \textbf{发布时间} & \textbf{2025年目标} & \textbf{2030年目标} & \textbf{状态} \\
\midrule
《"十四五"可再生能源发展规划》 & 2022.06 & 330GW（光伏） & — & 已落地 \\
《促进新时代新能源高质量发展实施方案》 & 2022.05 & — & 12亿千瓦（风+光） & 已落地 \\
《"十四五"新型储能发展实施方案》 & 2022.03 & 30GW+（储能） & — & 已落地 \\
"双碳"目标 & 2020.09 & — & 12亿千瓦以上（风+光） & 国家战略 \\
\midrule
\rowcolor{lightgrey}
\multicolumn{5}{l}{\textit{"十五五"前瞻}} \\
光伏装机（前瞻） & — & — & 1500GW+ & 预期中 \\
新型储能（前瞻） & — & — & 150GW+ & 预期中 \\
\bottomrule
\end{tabularx}
\sourcenote{国家能源局、发改委、公开信息整理}
\end{exhibitbox}

从实际装机情况看，光伏和储能均持续超预期。2024年中国光伏新增装机约230GW，2025年预计约230--250GW，已大幅超过"十四五"规划目标。储能装机同样爆发式增长，2025年累计装机预计约60--70GW，已是"十四五"30GW目标的两倍以上。

光伏用功率半导体市场方面，2025年全球市场规模约30--35亿美元，2030年预计达60--70亿美元，IGBT约占60\%，MOSFET约占20\%，SiC约占20\%且处于快速增长中。储能用功率半导体增长更快，2025年全球市场约15--20亿美元，2030年预计达50--60亿美元，CAGR约25--30\%，是SiC器件渗透最快的应用场景之一。

\subsection{AI算力政策：新增长极的政策支撑}

AI算力基础设施建设是功率半导体需求的新兴边际增量，受到政策的大力推动。"东数西算"工程、《算力基础设施高质量发展行动计划》、《人工智能+行动方案》等政策持续加码算力基础设施建设。

根据政策目标，2025年总算力规模超过300EFLOPS，智能算力占比达到35\%；2030年智能算力占比有望超过50\%，总规模超1000EFLOPS。AI服务器和数据中心建设拉动服务器电源和供电架构升级，带动高端PMIC、DrMOS、GaN/SiC等功率半导体需求增长，单机价值量较传统服务器提升3--6倍。

数据中心能效标准的提高（如PUE≤1.3要求）进一步推动高效率功率器件的应用。GaN在48V到负载点的转换中优势明显，SiC在800V高压直流转换中优势突出，电源架构从传统12V供电向48V母线、800V HVDC演进的过程中，功率半导体单机价值量持续跃升。

\begin{exhibitbox}[表 8-10：A股功率半导体公司新能源政策敏感度排序]
\centering
\small
\begin{tabularx}{\textwidth}{l c X c}
\toprule
\textbf{公司} & \textbf{新能源业务占比} & \textbf{主要受益领域} & \textbf{政策敏感度} \\
\midrule
天岳先进 & 约70--80\% & SiC衬底（新能源车、光伏、储能） & \textcolor{riskred}{很高} \\
三安光电 & 约30--35\% & SiC器件（新能源车、光伏、储能） & \textcolor{riskred}{很高} \\
斯达半导 & 约40--45\% & 新能源车IGBT模块、光伏逆变器 & \textcolor{riskred}{很高} \\
时代电气 & 约35--40\% & 新能源车电驱、储能PCS、光伏IGBT & \textcolor{riskred}{高} \\
扬杰科技 & 约30--35\% & 光伏、储能、新能源车 & \textcolor{riskred}{高} \\
华润微 & 约25--30\% & 光伏MOS/IGBT、储能、新能源车 & \textcolor{riskamber}{中高} \\
士兰微 & 约20--25\% & 新能源车IGBT、光伏 & \textcolor{riskamber}{中高} \\
新洁能 & 约25--30\% & 光伏MOSFET、储能 & \textcolor{riskamber}{中高} \\
东微半导 & 约20--25\% & 充电桩、光伏、储能 & \textcolor{riskamber}{中} \\
晶丰明源 & 约10--15\% & 储能、光伏电源管理 & \textcolor{riskamber}{低--中} \\
\bottomrule
\end{tabularx}
\sourcenote{各公司年报、IR交流、本研究测算}
\end{exhibitbox}

从政策敏感度看，SiC产业链企业（天岳先进、三安光电）和车规IGBT龙头（斯达半导）对新能源政策最为敏感，新能源业务占比高且行业处于快速成长期，政策边际变化对业绩和估值影响显著。传统功率IDM企业（时代电气、扬杰科技、华润微、士兰微）新能源业务占比适中，兼具成长和安全垫。

\begin{keyinsight}[新能源政策联动效应]
新能源汽车、光伏、储能、AI算力四大下游的政策支撑共同构成功率半导体需求的"四引擎"。各政策的时间维度互补——新能源车购置税明确至2027年、光伏储能"十五五"目标持续上调、AI算力建设处于加速期——为行业提供持续的需求催化。投资上，建议沿政策敏感度和业绩确定性两个维度配置，兼顾弹性与稳健性。
\end{keyinsight}

\section{地缘政治风险情景分析}

\subsection{情景分析框架}

我们从国内政策力度、出口管制强度、国际贸易环境三个核心维度构建地缘政治风险情景分析框架，对功率半导体行业面临的政策不确定性进行系统评估。三种情景的概率分布为：乐观情景20--25\%、中性情景50--55\%、悲观情景25--30\%。

\begin{exhibitbox}[表 8-11：三种情景核心假设对比]
\centering
\small
\begin{tabularx}{\textwidth}{X L{4.2cm} L{4.2cm} L{4.2cm}}
\toprule
\textbf{维度} & \textbf{乐观情景（25\%）} & \textbf{中性情景（50\%）} & \textbf{悲观情景（25\%）} \\
\midrule
国内政策 & 大幅加码：大基金超预期投放，"十五五"目标大幅上调 & 按部就班：政策按既定节奏推进，目标适度上调 & 不及预期：支持力度放缓，新能源补贴退坡 \\
出口管制 & 维持现状：管制未进一步升级，达成部分共识 & 小幅升级：在预期范围内，功率半导体受间接影响 & 大幅升级：成熟制程设备和SiC/GaN纳入管制 \\
国际贸易 & 环境改善：关税壁垒降低，CBAM不扩围 & 常态摩擦：关税维持，CBAM逐步推进 & 摩擦加剧：关税大幅提升，CBAM扩围至半导体 \\
\midrule
行业增速CAGR & 15--18\% & 10--12\% & 5--7\% \\
国产化提升 & 5--8\%/年 & 3--5\%/年 & 1--3\%/年 \\
板块表现 & +30\,\~{}\,+50\% & 0\,\~{}\,+20\% & -20\,\~{}\,-30\% \\
\bottomrule
\end{tabularx}
\sourcenote{本研究情景分析框架}
\end{exhibitbox}

\color{black}

\subsection{乐观情景：政策红利集中释放}

在乐观情景下，国内政策支持力度超预期，大基金三期投放速度加快且功率半导体和第三代半导体获重点支持，"十五五"规划新能源和算力目标大幅上调。出口管制未进一步升级，中美在科技领域达成部分共识。全球贸易环境改善，关税壁垒降低，CBAM暂不扩围至半导体。

此情景下，2030年国内功率半导体市场规模较基准情景高出15--20\%，各环节国产化率提升5--10个百分点，车规IGBT国产化率达25--30\%，SiC衬底达35--40\%。行业平均净利润率提升2--3个百分点，板块估值中枢上移15--25\%。

受益标的排序为：三安光电（SiC/GaN全产业链布局，政策受益最大）、天岳先进（SiC衬底龙头，国产替代加速）、斯达半导（车规IGBT龙头，需求爆发）、时代电气（IDM龙头，新能源+电网双受益）、华润微（功率IDM，AI电源布局领先）。

\subsection{中性情景：摩擦常态，稳步发展}

中性情景是我们认为概率最高的基准情景。国内政策按既定节奏推进，大基金三期正常投放，"十五五"规划符合预期。出口管制有所升级但在预期范围内，主要影响先进制程和AI芯片，功率半导体受间接影响。贸易摩擦常态化，关税水平维持，CBAM逐步推进但暂未扩围至半导体。

此情景下，国内功率半导体市场2025--2030年CAGR约10--12\%，国产化率年均提升3--5个百分点，车规IGBT国产化率达15--20\%，SiC衬底达25--30\%。行业平均净利润率基本稳定或小幅提升，板块震荡上行，估值中枢基本稳定但结构性分化加剧。

受益标的排序为：时代电气（估值洼地，安全边际高）、扬杰科技（业绩确定性强，财务质量优）、华润微（产能释放+产品结构升级）、斯达半导（车规龙头，增长确定性高）、士兰微（产品布局全面，AI+汽车双驱动）。

\subsection{悲观情景：管制升级，承压调整}

悲观情景下，出口管制大幅升级，成熟制程设备和材料纳入管制范围，SiC/GaN被明确纳入管制。国内政策支持力度不及预期，大基金投放节奏放缓，经济下行压力增大，新能源补贴退坡。贸易摩擦加剧，关税大幅提升，CBAM扩围至半导体，欧洲市场出口受阻。

此情景下，2030年国内市场规模较基准低10--15\%。短期受设备材料限制，扩产受阻，行业平均净利润率下降3--5个百分点，板块估值回调20--30\%，可能出现业绩下修+估值压缩的戴维斯双杀。但长期看，国产替代的大方向不变，只是节奏因外部压力短期放缓。

风险暴露从高到低排序为：三安光电（SiC/GaN受管制直接影响大，估值高）、天岳先进（SiC衬底，海外客户占比高，设备受限）、士兰微/华润微（扩产计划受设备限制影响大）、斯达半导（出口占比较高，估值较高）、时代电气（轨交安全垫，估值低，相对抗跌）。

相对抗跌标的包括时代电气（轨交业务提供约55\%收入安全垫，PE约20x）、扬杰科技（财务质量好，现金流健康，估值合理）、新洁能（Fabless模式轻资产，受设备管制影响小）。

\begin{keyinsight}[情景分析结论与投资启示]
综合三种情景，中性情景概率最高（50--55\%），行业处于政策红利与地缘风险的平衡状态，结构性机会大于系统性机会。投资策略上，建议以中性情景为基准配置，以确定性龙头为底仓，同时预留部分弹性仓位捕捉乐观情景下的超额收益。在悲观情景风险释放时，重点加仓被错杀的优质龙头。建议关注的核心跟踪指标包括：大基金投资进展、BIS管制规则更新、"十五五"规划政策、新能源装机数据、设备交付周期等。
\end{keyinsight}
